\documentclass[12pt]{article}
\usepackage[utf8]{inputenc}
\usepackage[T1]{fontenc}
\usepackage{hyperref}
\title{Spanish Sign Language Language Resources for K-12 Curriculum Development}
\author{Compiled by Claude Sonnet 5.5 (claude-sonnet-5-5) for Andr\'as Kornai}
\date{}
\begin{document}
\maketitle

\section{Introduction}
Spanish Sign Language (Lengua de signos espa\~nola, LSE; ISO 639-3 ssp) is the sign language of the Spanish deaf community. It was legally recognised by Ley 27/2007, of 23 October, which recognises the Spanish sign languages and regulates support for oral communication of deaf people. Most of what the search found is in Spanish and predates 2015, so current availability should be checked with the CNSE (Confederaci\'on Estatal de Personas Sordas).

\section{Primary Educational Resources}
The searches described the following CNSE educational materials. I did not see the materials themselves online, only descriptions of them.
\begin{itemize}
\item The CNSE portal ``Actividades en LSE'' for educators and sign-language specialists in schools, aimed at Infantil and Primaria. It includes DVD videos made with the Ministry of Education, curricular proposals and methodological guidance. Described at: \url{https://www.consumer.es/?p=233005}
\item The CNSE ``Proyecto ABC'' in a Murcia secondary school (2009), which placed sign-competent staff with deaf students: \url{https://www.cnse.es/media/k2/attachments/226.2009.pdf}
\item Other CNSE video materials named in the results (``En casa signamos'', ``En el cole signamos'') are described in the sources but no direct links were seen.
\end{itemize}

\section{Linguistic and Reference Materials}
\begin{itemize}
\item The CNSE Foundation's \emph{Diccionario normativo de la lengua de signos espa\~nola} (2008) is mentioned in the results, but no direct link was seen.
\item University of Granada course guide listing the online dictionary Spread The Sign, the DILSE mobile dictionary app and the Pinedo Peydr\'o (2007) dictionary: \url{https://grados.ugr.es/sites/grados/default/public/guias-firmadas/2024-2025/28611A6.pdf}
\item Academic articles on LSE teaching and deaf education: \url{https://dialnet.unirioja.es/descarga/articulo/3286682.pdf} and \url{https://cnlse.es/antiguos/15Gomez_yOtros.pdf}
\end{itemize}

\section{Religious and Community Resources}
The CNSE is the national deaf federation and the main community organisation in these results (\url{https://www.cnse.es/media/k2/attachments/226.2009.pdf}). A 2005 CCOO note citing the CNSE argues that sign language as a route into the curriculum can reduce school failure among deaf students: \url{https://hemerotecate.fe.ccoo.es/assets/20050616.pdf}. No religious resources were found.

\section{Digital and Audio Resources}
Spread The Sign and DILSE (above) are the digital dictionary resources seen in the results. No other video corpora were identified.

\section{Curriculum Development Context}
LSE has legal recognition and CNSE-developed school materials for Infantil and Primaria. Gaps seen in the results are secondary-level materials, consolidated support in schools, and verified current availability of the CNSE resources. A bilingual LSE--Spanish curriculum for deaf and hearing learners would build on these materials and on the CNSE's school experience.

\end{document}